\section{Discussion}
\label{sec:discussion}
\paragraph{What the video JEPA learned.} Trained on \val{pre.t1.clips} clips of one cell type
from one instrument, TiJepa-1 learned representations that classify microglia morphology
far better than hand-crafted features. Image models trained on millions of natural images
do as well, and a general-purpose video model used frozen and off the shelf, V-JEPA~2,
does better. On the label-free validation task TiJepa-1 is below raw pixels.
Future masking helps relative to the matched spatial-only ablation, modestly on morphology
and strongly on re-identification, but the predictor's apparent forecasting skill is
explained by the positional structure of its own target latents: it persists on clips in
which nothing moves. The representation is partially collapsed, and rank and downstream
performance move together across models. A practitioner who wants frozen features of
innate-cell video today should start from a large pretrained video model.

\paragraph{What public data can and cannot answer.} The coverage map makes the central gap
concrete: every public pairing of innate-cell imaging with a molecular or functional readout
is either at the same-cell level for states that are not memory (morphology classes, cell
type, metabolic lifetime, stimulus identity), or at the condition level with five conditions.
Where a simple measured reference exists (autofluorescence intensity for cell identity,
the photon image for lifetime) the learned features did not add to it. The memory states
themselves are measured only by restimulation assays that have never been imaged.

\paragraph{A constraint for the experiment that is missing.} The replicated observation
that $\beta$-glucan-trained human macrophages have less TNF transcript than untrained cells
4\,h after restimulation, in two datasets with different restimuli and donors, while the
same studies report more TNF-$\alpha$ protein at 24\,h, means that a single readout time can
show training as either enhanced or reduced. An imaging study of trained immunity therefore
needs its molecular ground truth at more than one time after restimulation, and its
imaging window should cover the hours in which the transcript and protein responses
diverge (Section~\ref{sec:future}).
